\documentclass{article}
%\usepackage[letterpaper,body={6.3in,9.15in},top=.8in,left=1.1in]{geometry}
\usepackage{fancyvrb,color,gretl}
\usepackage{amsmath}
\usepackage{graphicx}
\usepackage[authoryear]{natbib}
\usepackage[pdftex,hyperfootnotes=false]{hyperref}

\definecolor{steel}{rgb}{0.03,0.20,0.45}

\hypersetup{pdftitle={kmeans},
            pdfauthor={Artur Tarassow and Allin Cottrell},
            colorlinks=true,
            linkcolor=blue,
            urlcolor=red,
            citecolor=steel,
            bookmarksnumbered=true,
            plainpages=false
          }

\setlength{\parskip}{1ex}
\setlength{\parindent}{0pt}
\setcounter{secnumdepth}{1}
\setcounter{tocdepth}{1}

\renewcommand{\floatpagefraction}{.8}
\newcommand{\tightlist}{}

\newenvironment{funcdoc}
{\noindent\hrulefill\\[-12pt]}
{\medbreak}

\title{The \textsf{PairPlot} addon}
\author{Artur Tarassow and Allin Cottrell}
%\date

\begin{document}
\maketitle

This package produces grids of subplots depicting pairwise
relationships in a dataset.  This package maps each variable in a list
onto a column and row in a grid of multiple axes.

It can also represent an additional level of conditionalization with the
`factor' argument (a series) which plots different subsets of data in
different colors and/or point/circle types. This is used to resolve
elements on a third dimension.

Please report bugs or comments on the gretl mailing list, report an
issue on github ( \url{https://github.com/atecon/PairPlot/issues} ) or write
to \url{atecon@posteo.de}.

\section{Usage via the GUI}\label{usage-via-the-gui}

The function \texttt{PairPlot()} can be accessed through the GUI. The
menu entry is ``View-\textgreater Graph specified vars-\textgreater Pair
Plot''. However, via scripting the user can control many more details.

\section{Public function}\label{public-function}

\begin{verbatim}
PairPlot(const list L, const series factor[null], bundle opts[null])
\end{verbatim}

Main function for creating pair-plots.

\subsection{Parameters}\label{parameters}

\begin{itemize}
\tightlist
\item
  \texttt{L}: list, Series to plot against each other
\item
  \texttt{factor}: series, Categorical (numeric or string-valued)
  variable for mapping plot aspects to different colors and/or points/
  circles (optional, default: null). The factor series should have at
  most 16 distinct values. This series must be of type `discrete'. To
  set a series to discrete, use the command
  \texttt{setinfo\ \textless{}series\ name\textgreater{}\ -\/-discrete}.
\item
  \texttt{opts}: bundle, Optional for passing additional parameters. See
  below for explanation (default: null)
\end{itemize}

\subsection{Returns}\label{returns}

An integer of value `0' (FALSE) if no error occurs, otherwise 1 (TRUE).

\subsection{Options passed to the PairPlot() function by means of the
`opts'
bundle}\label{options-passed-to-the-pairplot-function-by-means-of-the-opts-bundle}

The user can control the following aspects by adding the respective
parameter to the `opts' bundle.

\begin{itemize}
\item
  \texttt{centroid}: string, Draw a symbol highlighting the centroid of
  the 2-dimensional data. Possible options are \texttt{mean},
  \texttt{median} and \texttt{none} (default: ``none''). If a
  \texttt{factor} series is passed, the centroid for each distinct
  factor is shown.
\item
  \texttt{centroid\_label}: bool, Show the label of the centroid
  statistics if TRUE, otherwise do not show if FALSE (default: FALSE).
\item
  \texttt{centroid\_linewidth}: scalar, Width of the line (default: 2).
\item
  \texttt{centroid\_pointsize}: scalar, Size of the point (default: 2).
\item
  \texttt{filename}: string, Path plus filename plus file type.
  (optional, default: \texttt{display} which draws the resulting plot
  directly at the screen). Currently figures can only be stored in
  \textbf{png} format.
\item
  \texttt{fontsize}: int, Control the font size of the labels (default:
  16) .
\item
  \texttt{grid}: bool, Draw a grid in the background if TRUE (=1)
  (default: FALSE)
\item
  \texttt{height}: scalar, Height of the canvas plot (default: 600, min.
  size: 100).
\item
  \texttt{title}: string: set an overall title for the plot(s).
\item
  \texttt{key}: bool, If the \texttt{factor} series is provided, a
  legend shows the color and point pattern for each distinct value of
  the \texttt{factor} variable. Default: 1 (TRUE).
\item
  \texttt{key\_fontsize} int, Control the font size for the key.
  Default: 14
\item
  \texttt{key\_position}: string, Controls the position of the legend in
  each subplot (use standard gnuplot options). default: ``top left''.
\item
  \texttt{offset\_level}: scalar, Provides a mechanism to put a boundary
  around the data inside of an auto-scaled graph. For details see:
  \url{https://gnuplot.sourceforge.net/docs\_6.0/loc14594.html} , Default:
  0.02.
\item
  \texttt{tmargin}: scalar, Top margin of the plot in lines. Positive
  values move the
\item
  \texttt{bmargin}: scalar, Bottom margin of the plot in lines.
\item
  \texttt{lmargin}: scalar, Left margin of the plot in lines.
\item
  \texttt{rmargin}: scalar, Right margin of the plot in lines.
\item
  \texttt{pointsize}: int, Control the size of points (default: 2). No
  support for circles.
\item
  \texttt{pointtype}: matrix, Control the type of points. Can be either
  an integer entry such that points associated to different levels of
  \texttt{factor} are of the same type. One may also provide a vector of
  integers for each distinct value of the \texttt{factor} series.
  Default: 7. See also: \url{https://i.imgur.com/xNZa7Rz.png}
\item
  \texttt{tics}: bool, Print x-tics and y-tics if 1 (TRUE, default), or
  not if 0 (FALSE).
\item
  \texttt{tics\_fontsize} int, Control the font size of the x-y tics.
  Default: 12
\item
  \texttt{transparency\_level}: int, Controls the transparency level of
  points and circles (between 0 (no transparency) and 255 (full
  transparency)). Default: 90.
\item
  \texttt{type}: string, Select the type of plot (default:
  ``triangle''). The following types are supported:
\item
  \textbf{(i)} ``matrix'': Scatterplot matrix which draws a scatter plot
  between all combinations of variables included in list \texttt{L}.
\item
  \textbf{(ii)} ``triangle'': Lower triangular version of the
  scatterplot ``matrix'' type.
\item
  \textbf{(iii)} ``row'': Plot the first series element in list
  \texttt{L} in a single row against all remaining list members.
\item
  \textbf{(iv)} ``column'': Works like type ``row'' but subplots are
  arranged in a single column instead.
\item
  \texttt{use\_circles}: bool, Draw circles instead of points if set to
  1 (TRUE), default: 0 (FALSE).
\item
  \texttt{width}: scalar, Width of the canvas plot (default: 900, min.
  size: 100).
\item
  \texttt{xlabel\_offset\_horizontal}: scalar, Horizontal offset for the
  X-axis label (default: 0). Negative value move the label to the left,
  and vice versa.
\item
  \texttt{xlabel\_offset\_vertical}: scalar, Vertical offset for the
  X-axis label (default: 0). Negative values move the label downwards,
  and vice versa.
\item
  \texttt{ylabel\_offset\_horizontal}: scalar, Horizontal offset for the
  Y-axis label (default: 0). Negative value move the label to the left,
  and vice versa.
\item
  \texttt{ylabel\_offset\_vertical}: scalar, Vertical offset for the
  Y-axis label (default: 0). Negative values move the label downwards,
  and vice versa.
\end{itemize}

\section{Changelog}\label{changelog}

\begin{itemize}
\tightlist
\item
  \textbf{October 2026}
  
  Switch help to \LaTeX\ (pdf) format.
  
\item
  \textbf{v1.1 (July 2025)}

  \begin{itemize}
  \tightlist
  \item
    For triangle-type of plot: improved x-/y-label positioning
  \item
    Make margins of the plot configurable via \texttt{tmargin},
    \texttt{bmargin}, \texttt{lmargin} and \texttt{rmargin} parameters.
  \item
    Added \texttt{xlabel\_offset\_horizontal},
    \texttt{xlabel\_offset\_vertical},
    \texttt{ylabel\_offset\_horizontal} and
    \texttt{ylabel\_offset\_vertical} parameter to control X-axis and
    Y-axis label offset.
  \end{itemize}
\item
  \textbf{v1.0 (February 2025)}

  \begin{itemize}
  \tightlist
  \item
    Bugfix: Catch error in case size of width or height is less than
    100.
  \item
    Bugfix: Sample 2 not working before as reference to dataset
    \texttt{iris.gdt} was missing.
  \item
    Ensure that the \texttt{factor} series is of type `discrete'. If
    not, an error message is thrown.
  \item
    Throw an error if the \texttt{factor} series has more than 16
    distinct values.
  \end{itemize}
\item
  \textbf{v0.99 (July 2024)}

  \begin{itemize}
  \tightlist
  \item
    Add new parameter \texttt{title} for setting an overall title
  \item
    Internal refactoring: Switch to gretl's built-in gridplot apparatus
    which means that all the graphics file formats supported by gretl
    can be used.
  \item
    It also means that usage via gretlcli in ``display'' mode works.
  \item
    In case the user's specification calls for just a single plot, no
    ``multiplot'' is created.
  \item
    The graphical interface is enhanced, with more options.
  \end{itemize}
\item
  \textbf{v0.98 (April 2024)}

  \begin{itemize}
  \tightlist
  \item
    Make width, height, font size and point size a function of the
    number of variables for better readability in case of many
    variables.
  \item
    Set parameter \texttt{transparency} per default to 90.
  \item
    Minor update of the sample script
  \end{itemize}
\item
  \textbf{v0.97 (February 2024)}

  \begin{itemize}
  \tightlist
  \item
    Error in case of CLI-mode if a plot is expected to appear on the
    screen (filename=``display'')
  \end{itemize}
\item
  \textbf{v0.96 (February 2024)}

  \begin{itemize}
  \tightlist
  \item
    Set all terminals to `noenhanced' mode to respect eventual
    underscores in variable names
  \item
    Allow the parameters `centroid\_pointsize', `centroid\_linewidth',
    `width' and `height' to be scalar values instead of integers.
  \end{itemize}
\item
  \textbf{v0.95 (December 2023)}

  \begin{itemize}
  \tightlist
  \item
    Bugfix: Path with blanks failed under Windows sometimes.
  \item
    Bugfix: An error occurred with negative values in some cases
  \item
    Bugfix: An error occurred of the name of the last variable in the
    list starts with `e' (or `E')
  \item
    Fix handling of transparency
  \item
    Add \texttt{grid} option for (optionally) drawing a grid in the
    background
  \item
    Add options \texttt{key\_fontsize} and \texttt{tics\_fontsize}.
  \item
    Add support for drawing centroids of 2-dimensional space.
  \item
    Remove the \texttt{font} option for setting the font (controlled by
    gretl now)
  \item
    Increase default font size to 12
  \item
    Switch to markdown-based help file
  \item
    Increase minimum gretl version to 2022d
  \end{itemize}
\item
  \textbf{v0.9 (October 2022)}

  \begin{itemize}
  \tightlist
  \item
    Fix bug with \texttt{const} inheritance: its illegal for a function
    which receives an argument A that is declared as const to pass A on
    in pointer form to another function which does not promise to treat
    it as \texttt{const}.
  \end{itemize}
\item
  \textbf{v0.8 (January 2022)}

  \begin{itemize}
  \tightlist
  \item
    Handle eventual missing values in passed series.
  \item
    Bugfix of gnuplot's \texttt{offset} option: More than at maximum 4
    were passed in case of more than 4 features to plot.
  \item
    Bugfix: Works now with more than 8 features.
  \end{itemize}
\item
  \textbf{v0.7 (October 2021)}

  \begin{itemize}
  \tightlist
  \item
    Enforce full sample instead of using \texttt{-\/-replace} option
  \item
    Changed ordering of the colors when plotting distinct factors
  \item
    Change the color palette to ``dark2''
  \end{itemize}
\item
  \textbf{v0.6 (March 2021)}

  \begin{itemize}
  \tightlist
  \item
    New API for \texttt{PairPlot()} -- backward-incompatible
  \item
    update help text
  \item
    enforce decimal point when print numerical data to gnuplot file
  \item
    substantial internal refactoring and clean up
  \item
    gain of speed
  \item
    new minimum gretl version is 2020e
  \end{itemize}
\item
  \textbf{v0.5 (Oct.~2019)}

  \begin{itemize}
  \tightlist
  \item
    Initial version
  \end{itemize}
\end{itemize}



\end{document}

%%% Local Variables:
%%% mode: latex
%%% TeX-master: t
%%% End:
